\chapter{Całki na zbiorach nieograniczonych i z osobliwościami}
\label{ma2:ch:improper-euclidean}

W analizie matematycznej I omówiono całki niewłaściwe
funkcji jednej zmiennej, w tym zbieżność bezwzględną,
warunkową i wartość główną Cauchy'ego
(\atlasref{AM1-31::ma1:ch:improper-integrals}{rozdział 31}).
W przestrzeni $\R^n$ można zbliżać się do nieskończoności
lub do zbioru osobliwości na wiele sposobów.
Punktem odniesienia będzie całka Lebesgue'a:
jeżeli funkcja jest całkowalna, wynik nie zależy od
sposobu wyczerpania dziedziny.
Nie będziemy zakładać twierdzenia Fubiniego
ani wzoru na zamianę zmiennych, które omówimy później.

\section{Wyczerpania zbiorów mierzalnych}

\begin{definicja}
\label{ma2:def:integral-exhaustion}
Niech $A\subseteq\R^n$ będzie mierzalny.
Wyczerpaniem mierzalnym zbioru $A$ nazywamy
ciąg mierzalnych zbiorów $A_j$ takich, że
\[
 A_j\subseteq A_{j+1}\subseteq A,\qquad
 \lambda_n(A_j)<\infty,\qquad
 A=\bigcup_{j\ge1}A_j.
\]
Mówimy o wyczerpaniu z dokładnością do zbioru zerowego,
jeżeli zamiast ostatniej równości zachodzi
$\lambda_n(A\setminus\bigcup_j A_j)=0$.
\end{definicja}

Każdy zbiór mierzalny ma wyczerpanie, na przykład
$A_j=A\cap[-j,j]^n$. Jeśli $f:A\to\R$ jest
mierzalna, możemy dodatkowo wymagać, aby $f$
była ograniczona na każdym $A_j$: wystarczy
przyjąć $A_j=A\cap[-j,j]^n\cap\{|f|\le j\}$.
Dla funkcji nieujemnej o wartościach rozszerzonych
taka modyfikacja pomija zbiór $\{f=+\infty\}$;
wolno go pominąć bez wpływu na całkę tylko wtedy,
gdy ma miarę zero.

\begin{twierdzenie}[Niezależność od wyczerpania]
\label{ma2:thm:improper-exhaustion}
Niech $(A_j)$ będzie wyczerpaniem mierzalnym
zbioru $A$ z dokładnością do zbioru zerowego.
\begin{enumerate}[label=(\roman*)]
\item Jeżeli $f:A\to[0,\infty]$ jest mierzalna, to
\[
 \int_A f(x)\,\mathrm d\lambda_n(x)
 =\lim_{j\to\infty}\int_{A_j}f(x)\,\mathrm d\lambda_n(x)
 \quad\text{w }[0,\infty].
\]
\item Jeżeli $f:A\to\R$ jest całkowalna, to
\[
 \int_A f(x)\,\mathrm d\lambda_n(x)
 =\lim_{j\to\infty}\int_{A_j}f(x)\,\mathrm d\lambda_n(x).
\]
Ponadto
$\int_{A\setminus A_j}|f(x)|\,\mathrm d\lambda_n(x)\to0$.
\item Dla mierzalnej $f:A\to\R$ warunek
\[
 \sup_j\int_{A_j}|f(x)|\,\mathrm d\lambda_n(x)<\infty
\]
jest równoważny całkowalności $f$ na $A$.
\end{enumerate}
\end{twierdzenie}
\begin{proof}
Niech $N=A\setminus\bigcup_j A_j$.
Przedłużenia ograniczeń $f|_{A_j}$ przez zero
na $A$ oznaczamy przez $f\chi_{A_j}$, przy
konwencji $0\cdot(+\infty)=0$.
Dla $f\ge0$ tworzą one ciąg niemalejący
o granicy $f$ poza $N$ i zero na $N$.
Twierdzenie o zbieżności monotonicznej
\ref{ma2:thm:lebesgue-monotone} oraz niezmienniczość
całki po zmianie na zbiorze zerowym
(\atlasref{ME-07::thm:ae-integral-invariance}{teoria miary i całki})
dają punkt (i).

W punkcie (ii) funkcje $f\chi_{A_j}$ są
zmajoryzowane przez $|f|$ i zbiegają do $f$
prawie wszędzie. Stosujemy twierdzenie
\ref{ma2:thm:lebesgue-dominated}.
Do funkcji $|f|\chi_{A\setminus A_j}$ stosujemy
to samo twierdzenie: zbiegają prawie wszędzie
do zera i mają tę samą całkowalną majorantę.
Punkt (iii) wynika z (i) zastosowanego do $|f|$,
gdyż całki nieujemne na $A_j$ są niemalejące,
a ich granica jest ich supremum.
\end{proof}

Punkty (i) i (ii) można stosować również do
niemonotonicznego ciągu zbiorów $A_j\subseteq A$,
jeżeli $\chi_{A_j}\to\chi_A$ prawie wszędzie.
Dla całkowalnej $f$ wynika to z
\ref{ma2:cor:lebesgue-moving-sets}.
Dla nieujemnej $f$ o nieskończonej całce
lemat Fatou \ref{ma2:thm:lebesgue-fatou} daje wówczas
$\liminf_j\int_{A_j}f\,\mathrm d\lambda_n=+\infty$,
gdyż funkcje przedłużone przez zero
zbiegają prawie wszędzie do $f$.
Dla skończonej całki nieujemnej stosujemy
zbieżność zmajoryzowaną po pominięciu zbioru,
na którym $f=+\infty$.

\section{Wyczerpania geometryczne obszaru}

Niech $U\subseteq\R^n$ będzie otwarty i
$f:U\to\R$ ciągła. Nawet jeśli $U$ jest
nieograniczony lub $f$ nie jest ograniczona
w pobliżu brzegu, na zwartych podzbiorach
$U$ funkcja $f$ jest ograniczona.

\begin{lemat}[Wyczerpanie kostkami]
\label{ma2:lem:improper-jordan-exhaustion}
Każdy otwarty zbiór $U\subseteq\R^n$ ma
wyczerpanie $(K_j)$ złożone ze zwartych
zbiorów mierzalnych w sensie Jordana,
zawartych w $U$.
\end{lemat}
\begin{proof}
Dla $j\ge1$ rozpatrujemy wszystkie domknięte
kostki siatki diadycznej o boku $2^{-j}$,
zawarte w $U\cap[-j,j]^n$.
Ich suma $K_j$ jest skończona, gdyż
$[-j,j]^n$ jest ograniczony. Przy braku
takich kostek przyjmujemy $K_j=\varnothing$.
Każdy $K_j$ jest zwarty. Brzeg skończonej
sumy domkniętych kostek zawiera się w sumie
ich brzegów: poza tą sumą każdy punkt jest
albo we wnętrzu jednej kostki, albo poza
wszystkimi kostkami w pewnym otoczeniu.
Każdy brzeg kostki jest zawarty w skończonej
sumie hiperpłaszczyzn współrzędnościowych,
które mają miarę zero
(\atlasref{ME-04::lem:hyperplane-null}{lemat o hiperpłaszczyznach}).
Z twierdzenia \ref{ma2:thm:jordan-criterion}
wynika mierzalność $K_j$ w sensie Jordana.

Każda kostka wybrana na etapie $j$ jest sumą
$2^n$ kostek siatki etapu $j+1$, również
zawartych w $U\cap[-(j+1),j+1]^n$.
Zatem $K_j\subseteq K_{j+1}$.
Dla każdego $x\in U$ istnieje $r>0$
takie, że kula euklidesowa $B(x,r)$
zawiera się w $U$. Wybieramy $j$
tak duże, by $j>\|x\|_\infty+r$
oraz $\sqrt n\,2^{-j}<r$.
Każda kostka tej siatki zawierająca $x$
ma średnicę $\sqrt n\,2^{-j}$,
a więc zawiera się w $B(x,r)$
i w $[-j,j]^n$. Stąd $x\in K_j$.
Otrzymujemy $\bigcup_j K_j=U$.
\end{proof}

\begin{wniosek}[Całka przez całki Riemanna na wyczerpaniu]
\label{ma2:cor:improper-riemann-exhaustion}
Niech $U\subseteq\R^n$ będzie otwarty,
$f:U\to\R$ ciągła, a $(K_j)$ dowolnym
rosnącym wyczerpaniem $U$ zwartymi zbiorami
mierzalnymi w sensie Jordana.
Jeżeli $f$ jest nieujemna lub całkowalna
w sensie Lebesgue'a, to
\[
 \int_U f(x)\,\mathrm d\lambda_n(x)
 =\lim_{j\to\infty} I_{K_j}(f|_{K_j}),
\]
gdzie $I_{K_j}$ oznacza całkę Riemanna
z rozdziału \ref{ma2:ch:lebesgue-euclidean}.
W przypadku $f\ge0$ dopuszczamy wartość $+\infty$.
Ponadto $f$ jest całkowalna wtedy i tylko wtedy,
gdy $\sup_j I_{K_j}(|f|\restriction_{K_j})<\infty$.
\end{wniosek}
\begin{proof}
Funkcje $f|_{K_j}$ i $|f|\restriction_{K_j}$ są ciągłe
i ograniczone. Twierdzenie
\ref{ma2:thm:lebesgue-jordan-agreement}
zapewnia istnienie całek Riemanna i ich
zgodność z całkami Lebesgue'a. Dla zbioru
pustego przyjmujemy całkę zero.
Stosujemy kolejno punkty (i)--(iii)
twierdzenia \ref{ma2:thm:improper-exhaustion}.
\end{proof}

Wniosek stosujemy np. na $U=\R^n\setminus S$,
gdzie $S$ jest domkniętym zbiorem osobliwości.
Jeżeli $S$ jest zerowy, wartości przypisane
funkcji na $S$ nie wpływają na całkę Lebesgue'a,
o ile rozpatrujemy mierzalne przedłużenie.
Nie oznacza to, że osobliwość jest nieszkodliwa:
o całkowalności decyduje zachowanie funkcji
w otoczeniu $S$.

\section{Osobliwości punktowe i zachowanie w nieskończoności}

Oznaczmy $\rho(x)=\|x\|_\infty=\max_i|x_i|$,
a przez $C_r=[-r,r]^n$ kostkę o promieniu $r>0$
w tej normie.

\begin{twierdzenie}[Kryterium potęgowe]
\label{ma2:thm:improper-power-test}
Niech $p\in\R$ i $n\ge1$.
\begin{enumerate}[label=(\roman*)]
\item Funkcja $\|x\|^{-p}$ jest całkowalna
na $\{x:0<\|x\|<1\}$ wtedy i tylko wtedy,
gdy $p<n$.
\item Funkcja $\|x\|^{-p}$ jest całkowalna
na $\{x:\|x\|\ge1\}$ wtedy i tylko wtedy,
gdy $p>n$.
\end{enumerate}
Twierdzenie zachodzi dla normy euklidesowej
i normy supremum.
\end{twierdzenie}
\begin{proof}
Najpierw rozpatrzmy normę supremum.
Dla $j\ge0$ połóżmy
\[
 S_j=C_{2^{-j}}\setminus C_{2^{-(j+1)}}.
\]
Zbiory są rozłączne i ich suma to
$C_1\setminus\{0\}$. Ich miary wynoszą
\[
 \lambda_n(S_j)
 =2^n(1-2^{-n})\,2^{-jn}.
\]
Na $S_j$ mamy $2^{-(j+1)}<\rho(x)\le2^{-j}$.
Dla dowolnego rzeczywistego $p$ stąd
\[
 2^{-|p|}2^{jp}\le \rho(x)^{-p}
 \le 2^{|p|}2^{jp}.
\]
Zatem, oznaczając $c_n=2^n(1-2^{-n})>0$,
\[
 c_n2^{-|p|}2^{j(p-n)}
 \le\int_{S_j}\rho(x)^{-p}\,\mathrm d\lambda_n(x)
 \le c_n2^{|p|}2^{j(p-n)}.
\]
Rozkład całki nieujemnej na rozłączne
zbiory (rozdział \ref{ma2:ch:lebesgue-euclidean})
i kryterium zbieżności szeregu geometrycznego
dają warunek $p<n$.
Brzeg $C_1$ ma miarę zero, więc jego
pominięcie nie zmienia wyniku.

W nieskończoności używamy rozłącznych
warstw $T_j=C_{2^{j+1}}\setminus C_{2^j}$,
$j\ge0$. Zachodzi
$\lambda_n(T_j)=c_n2^{(j+1)n}$, a
$2^j<\rho(x)\le2^{j+1}$.
Analogiczne oszacowania dają
\[
 c_n2^{n-|p|}2^{j(n-p)}
 \le\int_{T_j}\rho(x)^{-p}\,\mathrm d\lambda_n(x)
 \le c_n2^{n+|p|}2^{j(n-p)}.
\]
Suma jest skończona dokładnie wtedy,
gdy $p>n$.

Dla normy euklidesowej zachodzi
$\rho(x)\le\|x\|_2\le\sqrt n\,\rho(x)$,
więc dla $x\ne0$
\[
 n^{-|p|/2}\rho(x)^{-p}
 \le\|x\|_2^{-p}
 \le n^{|p|/2}\rho(x)^{-p}.
\]
Kula euklidesowa jednostkowa zawiera wnętrze
$C_{1/\sqrt n}$ i zawiera się w $C_1$;
na różnicy między tymi kostkami funkcje
są ograniczone i zbiór ma skończoną miarę.
Podobnie różnica między zewnętrzem kuli
a zewnętrzem $C_1$ zawiera się w ograniczonym
obszarze $\rho(x)\ge1/\sqrt n$.
W obu przypadkach dodatkowe części mają
skończone całki. Porównanie przenosi więc
oba kryteria na normę euklidesową.
\end{proof}

\begin{przyklad}[Osobliwość i zanik w różnych wymiarach]
\label{ma2:ex:improper-power-dimensions}
W $\R^3$ funkcja $1/\|x\|_2^2$ jest
całkowalna w przebitej kuli jednostkowej,
ale nie na jej zewnętrzu.
Funkcja $1/\|x\|_2^4$ ma własności odwrotne.
Przypadek graniczny $1/\|x\|_2^3$
nie jest całkowalny ani przy zerze, ani
w nieskończoności: każda z warstw
w dowodzie daje wkład ograniczony od dołu
przez tę samą dodatnią stałą.
Natomiast funkcja
\[
 f(x)=\frac{1}{\|x\|_2^2(1+\|x\|_2)^2},
 \qquad x\ne0,
\]
jest całkowalna na $\R^3\setminus\{0\}$:
dla $0<\|x\|_2\le1$ jest nie większa od
$\|x\|_2^{-2}$, a dla $\|x\|_2\ge1$
od $\|x\|_2^{-4}$. Obie części są całkowalne
na mocy twierdzenia \ref{ma2:thm:improper-power-test}.
\end{przyklad}

% Następny blok: dokładne całki radialne przez miarę kul,
% osobliwości wzdłuż podzbiorów i zależność granic
% niewłaściwych od wyczerpania; wartość główna Cauchy'ego.
